\chapter{Resultados e Discussão}
\label{cha:resultados}

% Fontes: tcc2/agent-reports/METRICAS_REPLAY.md e chapters/tabelas/metricas-replay.tex
% (scripts/metricas_replay.py); tcc2/agent-reports/BOOTSTRAP_REPLAY.md e
% chapters/tabelas/bootstrap-replay.tex (scripts/bootstrap_replay.py);
% tcc2/evidence/04-replay-abril-novembro/relatorios/evolucao_replay_v1_v2_v3_tcc.md
% (documento mestre, seções 1.2, 2 a 6); auditoria_posicoes_replay_2026-04_2026-11.md;
% replay_continuo_v2_2026-04_2026-11.md; replay_continuo_v3_2026-04_2026-11.md;
% tcc2/evidence/Resultados-comparados.md; tcc2/evidence/01-direcao-e-estado/ESTADO_ATUAL.md.
Este capítulo apresenta os resultados do Replay Contínuo Abril--Novembro de 2026, obtidos sob o desenho descrito no Capítulo~\ref{cha:metodologia}: seis carteiras --- a carteira Multiagentes e cinco referências determinísticas --- percorrem os mesmos 153 pregões, a partir da mesma cesta inicial, em três campanhas sucessivas da carteira Multiagentes (v1, v2 e v3). A leitura dos resultados obedece à distinção estabelecida na Seção~\ref{sec:met-hipoteses}. Recebem decisão inferencial apenas o contraste de H1, entre duas referências determinísticas, e a família de contrastes de H2 na versão v1, cujas regras foram fixadas antes do replay. As versões v2 e v3 foram desenhadas após a observação do mesmo período e são exploratórias: seus contrastes são calculados pelo mesmo procedimento, mas reportados apenas como descrição, e nenhum valor-$p$ obtido nelas rejeita H0. Todos os números referem-se a uma única trajetória de mercado e são, portanto, resultados dentro da amostra.

O capítulo está organizado em seis partes: o contexto de mercado do período (Seção~\ref{sec:res-mercado}); os resultados agregados das seis carteiras nas três versões (Seção~\ref{sec:res-agregados}); a decisão sobre as hipóteses (Seção~\ref{sec:res-hipoteses}); a evolução arquitetural entre as versões, que responde à pergunta secundária do estudo (Seção~\ref{sec:res-evolucao}); a discussão do alcance dos resultados (Seção~\ref{sec:res-discussao}); e, separadamente, uma observação preliminar do acompanhamento prospectivo, que não integra o replay (Seção~\ref{sec:obs-operacional}).

\section{Contexto de Mercado do Período}
\label{sec:res-mercado}

A Tabela~\ref{tab:res-mercado} compara, para cada ativo do universo, o preço de referência da cesta inicial (fechamento de 02/04/2026) com o fechamento do último pregão do replay (12/11/2026). Quatro dos cinco ativos caíram no período, com perdas expressivas em BBAS3 e VALE3; apenas WEGE3 encerrou acima do preço de entrada. As variações são somente de preço e não incluem proventos.

\begin{table}[htbp]
\centering
\caption{Variação de preço dos ativos do universo no período do Replay Contínuo.}
\label{tab:res-mercado}
\small
\begin{tabular}{lrrr}
\toprule
\textbf{Ativo} & \textbf{Referência de entrada, 02/04 (R\$)} & \textbf{Fechamento, 12/11 (R\$)} & \textbf{Variação (\%)} \\
\midrule
BBAS3 & 23{,}39 & 15{,}48 & $-33{,}82$ \\
VALE3 & 83{,}55 & 63{,}95 & $-23{,}46$ \\
PETR4 & 48{,}15 & 44{,}70 & $-7{,}17$ \\
ITUB4 & 43{,}30 & 41{,}68 & $-3{,}74$ \\
WEGE3 & 50{,}31 & 51{,}65 & $+2{,}66$ \\
\bottomrule
\end{tabular}
\par\smallskip\footnotesize Fonte: documento mestre do replay (\texttt{evolucao\_replay\_v1\_v2\_v3\_tcc.md}, Seção 1.2), a partir de \texttt{operacoes\_v3.csv} e dos relatórios de auditoria.
\end{table}

Como a cesta inicial imposta tem pesos próximos de 20\% por ativo (Seção~\ref{sec:met-desenho}), esse movimento define o ponto de partida comum: a carteira Buy \& Hold, que apenas mantém a cesta, perdeu 11,21\% no período. Além da tendência de queda, a trajetória inclui um recuo concentrado no fim de outubro, visível nas curvas de \textit{drawdown} das carteiras que permaneceram investidas (Figura~\ref{fig:replay-drawdown}). Esse regime de queda condiciona a comparação entre as carteiras, ponto retomado na Seção~\ref{sec:res-discussao}.

\section{Resultados Agregados das Carteiras}
\label{sec:res-agregados}

As Tabelas~\ref{tab:replay-risco} e~\ref{tab:replay-operacoes} reúnem as métricas da Seção~\ref{sec:met-metricas} para as três versões da carteira Multiagentes e para as cinco referências, que são idênticas nas três campanhas, conforme a auditoria de paridade (Seção~\ref{sec:arq-replay}). As Figuras~\ref{fig:replay-patrimonio} e~\ref{fig:replay-drawdown} mostram as curvas de patrimônio e de \textit{drawdown}.

% Gerado por tcc2/scripts/metricas_replay.py; não editar à mão.
\begin{table}[htbp]
\centering
\caption{Retorno e risco das carteiras no Replay Contínuo, 06/04/2026 a 12/11/2026 (153 pregões). Sharpe pela definição \texttt{H2\_SCIENTIFIC\_SHARPE\_DEFINITION\_V1} ($r_f = 0$); Sortino com $r_f = 0$ e MAR $= 0$.}
\label{tab:replay-risco}
\small
\setlength{\tabcolsep}{4pt}
\resizebox{\textwidth}{!}{%
\begin{tabular}{lrrrrrrc}
\toprule
Carteira & Retorno (\%) & Ret. anual. (\%) & Vol. anual. (\%) & Sharpe & Sortino & MDD (\%) & Pico--vale \\
\midrule
Multiagentes v1 & $-17{,}82$ & $-27{,}62$ & 23{,}17 & $-1{,}28$ & $-1{,}60$ & $-22{,}90$ & 14/04--12/08 \\
Multiagentes v2 & $-14{,}26$ & $-22{,}38$ & 19{,}33 & $-1{,}21$ & $-1{,}58$ & $-19{,}37$ & 14/04--06/11 \\
Multiagentes v3 & 0{,}07 & 0{,}11 & 16{,}01 & 0{,}09 & 0{,}13 & $-12{,}80$ & 14/04--10/06 \\
\midrule
Buy \& Hold & $-11{,}21$ & $-17{,}78$ & 22{,}08 & $-0{,}78$ & $-1{,}01$ & $-18{,}53$ & 05/10--05/11 \\
SMA 50/200 & $-11{,}13$ & $-17{,}67$ & 13{,}11 & $-1{,}42$ & $-1{,}77$ & $-16{,}61$ & 14/04--06/11 \\
Bollinger 20/2 & 4{,}89 & 8{,}18 & 11{,}77 & 0{,}73 & 1{,}16 & $-8{,}52$ & 30/04--05/06 \\
Equal Weight 63d & $-10{,}13$ & $-16{,}13$ & 21{,}61 & $-0{,}71$ & $-0{,}94$ & $-18{,}10$ & 05/10--05/11 \\
Mínima Variância 252/63 & $-4{,}19$ & $-6{,}80$ & 17{,}52 & $-0{,}32$ & $-0{,}47$ & $-14{,}95$ & 13/04--05/06 \\
\bottomrule
\end{tabular}}
\par\smallskip\footnotesize Fonte: \texttt{daily.csv} das campanhas \texttt{REPLAY-01} a \texttt{REPLAY-03}; 153 pregões, dentro da amostra.
\end{table}

\begin{table}[htbp]
\centering
\caption{Operações, custos e caixa das carteiras no Replay Contínuo. Giro = volume negociado estimado (custo / 0{,}03\%) dividido pelo patrimônio médio.}
\label{tab:replay-operacoes}
\small
\resizebox{\textwidth}{!}{%
\begin{tabular}{lrrrrrr}
\toprule
Carteira & Patrimônio final (R\$) & Operações (C/V) & Custos (R\$) & Giro & Caixa médio (\%) & Caixa final (\%) \\
\midrule
Multiagentes v1 & 82.178{,}71 & 86 (50/36) & 184{,}40 & 6{,}78 & 12{,}98 & 1{,}64 \\
Multiagentes v2 & 85.743{,}57 & 38 (22/16) & 119{,}87 & 4{,}34 & 24{,}69 & 45{,}86 \\
Multiagentes v3 & 100.067{,}36 & 81 (29/52) & 120{,}07 & 4{,}03 & 14{,}94 & 0{,}75 \\
\midrule
Buy \& Hold & 88.791{,}11 & 0 (0/0) & 0{,}00 & 0{,}00 & 0{,}70 & 1{,}60 \\
SMA 50/200 & 88.866{,}68 & 6 (2/4) & 32{,}23 & 1{,}16 & 31{,}39 & 39{,}29 \\
Bollinger 20/2 & 104.887{,}05 & 24 (12/12) & 151{,}26 & 4{,}92 & 49{,}06 & 1{,}44 \\
Equal Weight 63d & 89.869{,}29 & 15 (9/6) & 5{,}11 & 0{,}18 & 0{,}70 & 1{,}58 \\
Mínima Variância 252/63 & 95.814{,}84 & 14 (8/6) & 27{,}94 & 0{,}97 & 0{,}78 & 1{,}65 \\
\bottomrule
\end{tabular}}
\par\smallskip\footnotesize Fonte: \texttt{daily.csv} das campanhas \texttt{REPLAY-01} a \texttt{REPLAY-03}; 153 pregões, dentro da amostra.
\end{table}


% Gerado por tcc2/scripts/metricas_replay.py; requer \usepackage{pgfplots}.
\begin{figure}[htbp]
\centering
\begin{tikzpicture}
\begin{axis}[
  width=\textwidth, height=7.5cm, xmin=0, xmax=153,
  xtick={0,19,39,60,83,104,125,146},
  xticklabels={abr,mai,jun,jul,ago,set,out,nov},
  xlabel={Pregão (2026)}, ylabel={Patrimônio (R\$ mil)},
  ymajorgrids, grid style={gray!20}, tick label style={font=\footnotesize},
  label style={font=\small}, scaled y ticks=false,
  yticklabel style={/pgf/number format/use comma, /pgf/number format/1000 sep={.}},
  legend style={at={(0.5,-0.2)}, anchor=north, legend columns=3,
    font=\footnotesize, draw=none, /tikz/every even column/.append style={column sep=6pt}},
  legend cell align=left,
]
\addplot[black, thick, densely dotted] table[x=idx, y=eq_ma_v1, col sep=comma] {figures/replay-series.csv};
\addlegendentry{Multiagentes v1}
\addplot[black, thick, dashed] table[x=idx, y=eq_ma_v2, col sep=comma] {figures/replay-series.csv};
\addlegendentry{Multiagentes v2}
\addplot[black, very thick] table[x=idx, y=eq_ma_v3, col sep=comma] {figures/replay-series.csv};
\addlegendentry{Multiagentes v3}
\addplot[gray, thin] table[x=idx, y=eq_bh, col sep=comma] {figures/replay-series.csv};
\addlegendentry{Buy \& Hold}
\addplot[gray, dashdotted] table[x=idx, y=eq_sma, col sep=comma] {figures/replay-series.csv};
\addlegendentry{SMA 50/200}
\addplot[gray, mark=square*, mark repeat=15, mark size=1.4pt] table[x=idx, y=eq_boll, col sep=comma] {figures/replay-series.csv};
\addlegendentry{Bollinger 20/2}
\addplot[gray, densely dashed, mark=o, mark repeat=15, mark size=1.6pt, mark options={solid}] table[x=idx, y=eq_ew, col sep=comma] {figures/replay-series.csv};
\addlegendentry{Equal Weight 63d}
\addplot[gray, mark=triangle*, mark repeat=15, mark size=1.8pt] table[x=idx, y=eq_mv, col sep=comma] {figures/replay-series.csv};
\addlegendentry{Mínima Variância 252/63}
\end{axis}
\end{tikzpicture}
\caption{Patrimônio de fechamento das carteiras no Replay Contínuo, de 02/04/2026 (R\$ 100 mil) a 12/11/2026. As referências são idênticas nas três campanhas.}
\label{fig:replay-patrimonio}
\end{figure}


% Gerado por tcc2/scripts/metricas_replay.py; requer \usepackage{pgfplots}.
\begin{figure}[htbp]
\centering
\begin{tikzpicture}
\begin{axis}[
  width=\textwidth, height=7.5cm, xmin=0, xmax=153,
  xtick={0,19,39,60,83,104,125,146},
  xticklabels={abr,mai,jun,jul,ago,set,out,nov},
  xlabel={Pregão (2026)}, ylabel={Drawdown (\%)},
  ymajorgrids, grid style={gray!20}, tick label style={font=\footnotesize},
  label style={font=\small}, scaled y ticks=false,
  yticklabel style={/pgf/number format/use comma, /pgf/number format/1000 sep={.}},
  legend style={at={(0.5,-0.2)}, anchor=north, legend columns=3,
    font=\footnotesize, draw=none, /tikz/every even column/.append style={column sep=6pt}},
  legend cell align=left,
]
\addplot[black, thick, densely dotted] table[x=idx, y=dd_ma_v1, col sep=comma] {figures/replay-series.csv};
\addlegendentry{Multiagentes v1}
\addplot[black, thick, dashed] table[x=idx, y=dd_ma_v2, col sep=comma] {figures/replay-series.csv};
\addlegendentry{Multiagentes v2}
\addplot[black, very thick] table[x=idx, y=dd_ma_v3, col sep=comma] {figures/replay-series.csv};
\addlegendentry{Multiagentes v3}
\addplot[gray, thin] table[x=idx, y=dd_bh, col sep=comma] {figures/replay-series.csv};
\addlegendentry{Buy \& Hold}
\addplot[gray, dashdotted] table[x=idx, y=dd_sma, col sep=comma] {figures/replay-series.csv};
\addlegendentry{SMA 50/200}
\addplot[gray, mark=square*, mark repeat=15, mark size=1.4pt] table[x=idx, y=dd_boll, col sep=comma] {figures/replay-series.csv};
\addlegendentry{Bollinger 20/2}
\addplot[gray, densely dashed, mark=o, mark repeat=15, mark size=1.6pt, mark options={solid}] table[x=idx, y=dd_ew, col sep=comma] {figures/replay-series.csv};
\addlegendentry{Equal Weight 63d}
\addplot[gray, mark=triangle*, mark repeat=15, mark size=1.8pt] table[x=idx, y=dd_mv, col sep=comma] {figures/replay-series.csv};
\addlegendentry{Mínima Variância 252/63}
\end{axis}
\end{tikzpicture}
\caption{Drawdown (queda em relação ao pico anterior) das carteiras no Replay Contínuo.}
\label{fig:replay-drawdown}
\end{figure}


\paragraph{Patrimônio e retorno.} Ordenadas pelo patrimônio final, as carteiras ficaram assim: Bollinger 20/2 (R\$ 104.887,05; +4,89\%), Multiagentes v3 (R\$ 100.067,36; +0,07\%), Mínima Variância 252/63 (R\$ 95.814,84; $-4{,}19\%$), Equal Weight 63d (R\$ 89.869,29; $-10{,}13\%$), SMA 50/200 (R\$ 88.866,68; $-11{,}13\%$), Buy \& Hold (R\$ 88.791,11; $-11{,}21\%$), Multiagentes v2 (R\$ 85.743,57; $-14{,}26\%$) e Multiagentes v1 (R\$ 82.178,71; $-17{,}82\%$). As versões v1 e v2 terminaram abaixo das cinco referências; a v3 terminou acima de quatro delas e R\$ 4.819,69 abaixo da Bollinger 20/2. Apenas duas carteiras tiveram retorno positivo, ambas com perfil de compra em fraqueza e venda em força (Seção~\ref{sec:res-v3}).

\paragraph{Sharpe e Sortino.} Somente a Bollinger 20/2 (0,73) e a v3 (0,09) tiveram Sharpe positivo. O erro-padrão aproximado de um Sharpe anualizado estimado com 153 retornos diários fica em torno de 1,3 para todas as carteiras \citep{Lo2002}, de modo que diferenças dessa ordem não se distinguem de ruído amostral nesta janela. Há ainda uma propriedade do índice que afeta a ordenação quando o retorno médio é negativo: como a média negativa é dividida pela volatilidade, uma carteira menos volátil recebe um Sharpe mais negativo para o mesmo retorno. É o caso da SMA 50/200, cujo retorno ($-11{,}13\%$) praticamente coincide com o do Buy \& Hold ($-11{,}21\%$), com volatilidade (13,11\% contra 22,08\%) e \textit{drawdown} máximo ($-16{,}61\%$ contra $-18{,}53\%$) menores, mas cujo Sharpe ($-1{,}42$) é o mais baixo entre as oito carteiras. O Sortino reproduz a mesma ordenação. Essa propriedade é retomada na decisão sobre H1 e H2 (Seção~\ref{sec:res-hipoteses}).

\paragraph{\textit{Drawdown}.} A v1 teve o maior \textit{drawdown} máximo ($-22{,}90\%$, de 14/04 a 12/08), seguida da v2 ($-19{,}37\%$); a v3 ($-12{,}80\%$) ficou atrás apenas da Bollinger 20/2 ($-8{,}52\%$). Para Buy \& Hold e Equal Weight 63d, o máximo ocorreu entre 05/10 e 05/11, no recuo do fim do período. A v1, a v2 e a SMA 50/200 permaneceram 146 pregões consecutivos abaixo do pico anterior, isto é, praticamente todo o período a partir de meados de abril.

\paragraph{Operações, custos e caixa.} A v1 foi a carteira mais ativa: 86 operações posteriores à entrada, R\$ 184,40 de custos e giro de 6,78 vezes o patrimônio médio. A v2 reduziu as operações para 38 e a v3 realizou 81, com predominância de vendas (52 contra 29 compras). Entre as referências, a Bollinger 20/2 operou 24 vezes, com giro de 4,92, e a SMA 50/200 apenas 6, com quatro saídas totais. Os custos acumulados, no máximo R\$ 184,40, são pequenos diante das diferenças de patrimônio entre as carteiras, da ordem de milhares de reais; as diferenças decorrem, portanto, das decisões, e não do custo proporcional de 0,03\%, que não inclui custo fixo nem lote mínimo (Seção~\ref{sec:met-ameacas}). O caixa distingue as estratégias: a Bollinger 20/2 manteve em média 49,06\% do patrimônio em caixa, a SMA 50/200 terminou com 39,29\% e a v2, com 45,86\%, enquanto Buy \& Hold, Equal Weight e Mínima Variância permaneceram investidas.

\paragraph{Pregões sem análise.} A v1 concluiu os 153 pregões com análise completa. A v2 teve um pregão \texttt{NO\_ANALYSIS} (26/06), por falha de transporte em uma chamada de risco, e a v3 teve quatro (01/07, 01/09, 01/10 e 02/10), por respostas do gestor recusadas na validação determinística. Nesses pregões as posições foram preservadas e marcadas normalmente, sem decisão fabricada (Seção~\ref{sec:arq-replay}).

\section{Decisão sobre as Hipóteses}
\label{sec:res-hipoteses}

A Tabela~\ref{tab:replay-bootstrap} apresenta os 16 contrastes de Sharpe calculados pelo Bootstrap A (Seção~\ref{sec:met-inferencia}): o contraste de H1, os cinco contrastes da família principal de H2 (v1) e os dez contrastes descritivos das versões v2 e v3. Todos foram estimáveis, nenhum teve reamostra degenerada e todos usaram a mesma sequência de índices. A coluna ``Critério?'' indica se a regra numérica de superioridade foi atendida em cada linha; ela equivale a rejeitar H0 apenas nas linhas de H1 e de H2 (v1), enquanto nas linhas de v2 e v3 é descritiva, sem interpretação de teste (Seção~\ref{sec:met-inferencia}).
% Gerado por tcc2/scripts/bootstrap_replay.py; não editar à mão.
\begin{table}[htbp]
\centering
\caption{Comparações pareadas de Sharpe no Replay Contínuo (153 pregões) pelo \textit{bootstrap} A: estacionário, \textit{basic} centrado, $B=5000$, semente 20261008, bloco médio 10. $\Delta$ = Sharpe(A) $-$ Sharpe(B), $r_f = 0$. $p_+$ e $L_{95}$: teste unilateral de superioridade (regra congelada); $U_{95}$ e $p_-$: direção inversa, exploratória. Holm dentro de cada família de cinco referências.}
\label{tab:replay-bootstrap}
\setlength{\tabcolsep}{4pt}
\resizebox{\textwidth}{!}{%
\begin{tabular}{llrrrrrrc}
\toprule
Hipótese & A vs B & $\Delta$ & $L_{95}$ & $U_{95}$ & $p_+$ & $p_+$ Holm & $p_-$ & Critério? \\
\midrule
H1 & SMA 50/200 vs Equal Weight 63d & $-0{,}71$ & $-1{,}75$ & 0{,}31 & 0{,}881 & 0{,}881 & 0{,}119 & não \\
\midrule
H2 (v1) & MA v1 vs Buy \& Hold & $-0{,}50$ & $-1{,}76$ & 0{,}93 & 0{,}728 & 1{,}000 & 0{,}272 & não \\
H2 (v1) & MA v1 vs SMA 50/200 & 0{,}14 & $-1{,}32$ & 1{,}76 & 0{,}432 & 1{,}000 & 0{,}568 & não \\
H2 (v1) & MA v1 vs Bollinger 20/2 & $-2{,}00$ & $-4{,}47$ & 0{,}41 & 0{,}917 & 1{,}000 & 0{,}084 & não \\
H2 (v1) & MA v1 vs Equal Weight 63d & $-0{,}57$ & $-1{,}88$ & 0{,}88 & 0{,}753 & 1{,}000 & 0{,}248 & não \\
H2 (v1) & MA v1 vs Mínima Variância 252/63 & $-0{,}96$ & $-2{,}27$ & 0{,}29 & 0{,}899 & 1{,}000 & 0{,}102 & não \\
\midrule
H2 (v2) & MA v2 vs Buy \& Hold & $-0{,}44$ & $-1{,}34$ & 0{,}83 & 0{,}714 & 1{,}000 & 0{,}286 & não \\
H2 (v2) & MA v2 vs SMA 50/200 & 0{,}20 & $-1{,}02$ & 1{,}69 & 0{,}347 & 1{,}000 & 0{,}653 & não \\
H2 (v2) & MA v2 vs Bollinger 20/2 & $-1{,}94$ & $-4{,}27$ & 0{,}31 & 0{,}917 & 1{,}000 & 0{,}083 & não \\
H2 (v2) & MA v2 vs Equal Weight 63d & $-0{,}51$ & $-1{,}46$ & 0{,}77 & 0{,}747 & 1{,}000 & 0{,}253 & não \\
H2 (v2) & MA v2 vs Mínima Variância 252/63 & $-0{,}90$ & $-1{,}96$ & 0{,}22 & 0{,}912 & 1{,}000 & 0{,}088 & não \\
\midrule
H2 (v3) & MA v3 vs Buy \& Hold & 0{,}86 & $-0{,}30$ & 2{,}13 & 0{,}108 & 0{,}432 & 0{,}892 & não \\
H2 (v3) & MA v3 vs SMA 50/200 & 1{,}50 & 0{,}25 & 3{,}02 & 0{,}024 & 0{,}121 & 0{,}976 & sim$^{*}$ \\
H2 (v3) & MA v3 vs Bollinger 20/2 & $-0{,}64$ & $-1{,}79$ & 0{,}59 & 0{,}812 & 0{,}812 & 0{,}188 & não \\
H2 (v3) & MA v3 vs Equal Weight 63d & 0{,}79 & $-0{,}34$ & 2{,}00 & 0{,}122 & 0{,}432 & 0{,}878 & não \\
H2 (v3) & MA v3 vs Mínima Variância 252/63 & 0{,}40 & $-1{,}09$ & 1{,}77 & 0{,}334 & 0{,}668 & 0{,}666 & não \\
\bottomrule
\end{tabular}}
\par\smallskip\footnotesize v1: execução pré-especificada; v2 e v3: iterações exploratórias desenhadas após observar o mesmo período (p-valores otimistas). Critério: $\hat{\Delta} > 0$, $p_+ < 0{,}05$ e $L_{95} > 0$; só rejeita $H_0$ em H1 e H2 (v1). $^{*}$Apenas descritivo: v2 e v3 não recebem decisão inferencial (Seção~\ref{sec:met-inferencia}); não se mantém após Holm. Fonte: \texttt{daily.csv} das campanhas \texttt{REPLAY-01} a \texttt{REPLAY-03}; \texttt{src/experiments/h2\_evaluation.py}.
\end{table}


\subsection{H1: tendência contra pesos iguais}
\label{sec:res-h1}

O contraste entre SMA 50/200 e Equal Weight 63d resultou em $\hat{\Delta} = -0{,}71$, com $p_{+} = 0{,}881$ e $L_{95} = -1{,}75$. Como a estimativa é negativa, o critério de superioridade não é atendido e \textbf{H0 não é rejeitada para H1}: nesta trajetória, a regra de tendência não apresentou Sharpe líquido superior ao da carteira de pesos iguais. O intervalo bilateral de 90\%, $[-1{,}75;\ 0{,}31]$, contém o zero, e o valor-$p$ exploratório de inferioridade ($p_{-} = 0{,}119$) também não indica desempenho inferior.

O sinal negativo da estimativa deve ser lido com a propriedade descrita na Seção~\ref{sec:res-agregados}. A SMA 50/200 realizou seis operações, zerou quatro posições e terminou com 39,29\% do patrimônio em caixa; seu retorno ($-11{,}13\%$) ficou 1,00 ponto percentual abaixo do da Equal Weight 63d ($-10{,}13\%$), mas com volatilidade e \textit{drawdown} máximo menores. Com retornos médios negativos nas duas carteiras, a menor volatilidade da SMA aprofundou seu Sharpe negativo, de modo que a diferença de $-0{,}71$ reflete mais essa propriedade do índice do que uma diferença de retorno. Por envolver duas carteiras determinísticas, calculadas antes de qualquer inferência paga, o resultado de H1 não é afetado pela possível contaminação do modelo nem pelo ajuste das versões v2 e v3. Seu alcance é restrito a uma regra de tendência com uma única parametrização, em um período de queda generalizada (Seção~\ref{sec:met-ameacas}).

\subsection{H2 na versão pré-especificada}
\label{sec:res-h2}

Os cinco contrastes da v1 foram negativos em quatro casos --- Buy \& Hold ($\hat{\Delta} = -0{,}50$; $p_{+} = 0{,}728$), Bollinger 20/2 ($-2{,}00$; $0{,}917$), Equal Weight 63d ($-0{,}57$; $0{,}753$) e Mínima Variância 252/63 ($-0{,}96$; $0{,}899$) --- e positivo apenas contra a SMA 50/200 ($+0{,}14$; $p_{+} = 0{,}432$). O enunciado de H2, de superioridade em relação ao conjunto das referências, exige que os cinco contrastes atendam ao critério (teste de interseção--união); nenhum o atende. Com o ajuste de Holm, todos os valores $p_{+}$ da família passam a 1,000, o que também afasta a superioridade em relação a referências individuais. Assim, \textbf{H0 não é rejeitada para H2 na versão v1}, e a hipótese de que o sistema multiagente apresente desempenho ajustado ao risco superior aos \textit{benchmarks} não é sustentada no replay.

O único contraste positivo não corresponde a desempenho melhor: a v1 terminou 6,69 pontos percentuais abaixo da SMA 50/200 em retorno, e a estimativa positiva decorre do Sharpe particularmente baixo da referência. Os critérios complementares, reportados de forma descritiva, apontam na mesma direção: a v1 teve o maior \textit{drawdown} máximo ($-22{,}90\%$) e o maior giro (6,78) entre as oito carteiras, e Sortino ($-1{,}60$) superior apenas ao da SMA 50/200. Na direção inversa, exploratória, nenhum contraste indica inferioridade a 5\%; o menor valor, $p_{-} = 0{,}084$, ocorre contra a Bollinger 20/2, com $\hat{\Delta} = -2{,}00$ e intervalo de 90\% de $[-4{,}47;\ 0{,}41]$. Uma diferença de dois pontos de Sharpe e de R\$ 22.708,34 em patrimônio final não é, portanto, distinguível de ruído amostral ao nível de 5\% nesta janela, o que ilustra o baixo poder discutido na Seção~\ref{sec:met-inferencia}.

\subsection{Contrastes descritivos das versões exploratórias}
\label{sec:res-h2-exploratorio}

A v2 repete o padrão da v1: contrastes negativos contra Buy \& Hold, Bollinger 20/2, Equal Weight 63d e Mínima Variância 252/63, positivo apenas contra a SMA 50/200 ($+0{,}20$; $p_{+} = 0{,}347$), sem nenhum $p_{+}$ inferior a 0,05. A v3 apresenta contrastes positivos contra Buy \& Hold ($+0{,}86$; $p_{+} = 0{,}108$), SMA 50/200 ($+1{,}50$; $p_{+} = 0{,}024$), Equal Weight 63d ($+0{,}79$; $p_{+} = 0{,}122$) e Mínima Variância 252/63 ($+0{,}40$; $p_{+} = 0{,}334$), e negativo contra a Bollinger 20/2 ($-0{,}64$; $p_{+} = 0{,}812$).

No contraste entre v3 e SMA 50/200, o valor $p_{+} = 0{,}024$ e o limite $L_{95} = 0{,}25$ atendem numericamente ao critério de superioridade. Conforme a Seção~\ref{sec:met-inferencia}, esse resultado é descritivo: as regras da v3 foram escolhidas com conhecimento dos mesmos 153 pregões e do desempenho das referências, e nenhum ajuste de multiplicidade corrige a busca implícita que produziu a versão \citep{White2000, Bailey2017}. Um valor-$p$ inferior a 0,05 nessa condição não tem interpretação de teste e não rejeita H0. Três observações reforçam essa leitura: mesmo dentro da família da v3, o valor ajustado por Holm é 0,121; a referência envolvida é a de menor Sharpe do período, rebaixado pela propriedade discutida na Seção~\ref{sec:res-h1}; e, contra a melhor referência, a Bollinger 20/2, o contraste da v3 é negativo.

\subsection{Síntese sobre H0}
\label{sec:res-h0}

Nas duas famílias que recebem decisão inferencial, H0 não é rejeitada: nem a regra de tendência superou a carteira de pesos iguais (H1), nem a carteira Multiagentes pré-especificada superou as referências (H2, v1). A não rejeição de H0 não demonstra igualdade de desempenho. Os intervalos de 90\% da Tabela~\ref{tab:replay-bootstrap} têm amplitude de 2,1 a 4,9 pontos de Sharpe, e o erro-padrão aproximado do Sharpe de cada carteira é da ordem de 1,3; nesse cenário, os dados são compatíveis tanto com a ausência de diferença quanto com diferenças economicamente relevantes em qualquer direção. Além disso, a inferência é condicional à única trajetória observada do modelo: com uma execução por versão e temperatura 1,0, a variabilidade entre execuções não está representada nos intervalos (Seção~\ref{sec:met-inferencia}). Sortino, \textit{drawdown} e giro permanecem critérios descritivos, sem contraste inferencial.

\section{Evolução Arquitetural entre as Versões}
\label{sec:res-evolucao}

A pergunta secundária do estudo (Seção~\ref{sec:met-caracterizacao}) é arquitetural: mantidos o modelo, a quantidade de agentes, os dados e a execução, quanto do comportamento da carteira Multiagentes depende da especialização dos analistas e das regras determinísticas aplicadas sobre a proposta do gestor. As intervenções de cada versão estão descritas na Seção~\ref{sec:arq-evolucao} (Tabelas~\ref{tab:evolucao-versoes} e~\ref{tab:constantes-v3}); esta seção apresenta seus efeitos observados. Como as cinco referências são idênticas nas três campanhas, as diferenças entre versões pertencem à carteira Multiagentes, mas são atribuíveis ao conjunto de intervenções de cada versão, e não a mecanismos isolados. A Tabela~\ref{tab:res-evolucao} resume os indicadores comparados.

\begin{table}[htbp]
\centering
\caption{Indicadores de comportamento da carteira Multiagentes nas três versões do Replay Contínuo.}
\label{tab:res-evolucao}
\small
\begin{tabular}{>{\raggedright\arraybackslash}p{6.6cm} rrr}
\toprule
\textbf{Indicador} & \textbf{v1} & \textbf{v2} & \textbf{v3} \\
\midrule
Patrimônio final (R\$) & 82.178{,}71 & 85.743{,}57 & 100.067{,}36 \\
Posição entre as seis carteiras & 6.ª & 6.ª & 2.ª \\
Operações pós-entrada (compras/vendas) & 86 (50/36) & 38 (22/16) & 81 (29/52) \\
Saídas totais (posição zerada) & 9 & 10 & 0 \\
Pareceres unânimes 5--0 (\%) & 83{,}01 & 15{,}13 & 4{,}03 \\
Pareceres com consenso mínimo 3--2 (\%) & 7{,}84 & 65{,}66 & 45{,}10 \\
RSI médio do sinal COMPRA / VENDA & 63{,}17 / 35{,}07 & 57{,}42 / 38{,}67 & 35{,}11 / 68{,}96 \\
RSI médio das compras / vendas executadas & 60{,}51 / 36{,}06 & 58{,}54 / 40{,}98 & 39{,}26 / 69{,}00 \\
Ativos com preço médio de venda $>$ compra & 1 de 5 & 1 de 5 & 5 de 5 \\
Pregões com 3 ou mais ativos de peso $\geq$ 2\% & 52 & 99 & 153 \\
Peso máximo de um ativo no fechamento (\%) & 77{,}49 & 32{,}23 & 35{,}87 \\
Caixa médio / final (\% do patrimônio) & 12{,}98 / 1{,}64 & 24{,}69 / 45{,}86 & 14{,}94 / 0{,}75 \\
Eventos de proventos recebidos (de 20) & 9 & 9 & 20 \\
Ajustes registrados pelo compilador & 0 & 26 & 80 \\
Pregões sem análise (\texttt{NO\_ANALYSIS}) & 0 & 1 & 4 \\
Custo conhecido estimado de inferência (US\$) & 7{,}44 & 8{,}92 & 8{,}63 \\
\bottomrule
\end{tabular}
\par\smallskip\footnotesize Fonte: documento mestre do replay (\texttt{evolucao\_replay\_v1\_v2\_v3\_tcc.md}, Resumo e Seções 5.1 a 5.9); pareceres com análise completa: 765 (v1), 760 (v2) e 745 (v3).
\end{table}

\subsection{Diagnóstico da versão v1}
\label{sec:res-v1}

A auditoria da campanha v1 descartou a hipótese de ativo esquecido --- os cinco ativos foram analisados em todos os pregões, em 765 pareceres ativo-pregão --- e isolou cinco achados empíricos.

\begin{enumerate}
    \item \textbf{Trava de caixa.} O caixa livre ficou abaixo de R\$ 200 em 69 dos 153 pregões, e o compilador estrito proibia reduzir uma posição cuja recomendação não fosse VENDA. A VALE3, zerada em 01/07, permaneceu fora da carteira por 94 pregões, até o fim da campanha, apesar de 14 recomendações de COMPRA aprovadas pelo risco entre 25/08 e 14/09, todas respondidas com HOLD; a ITUB4 recebeu 27 recomendações consecutivas de COMPRA (03/09 a 16/10) sem caixa para executá-las.
    \item \textbf{Compra no topo e venda no fundo.} O preço médio das compras superou o das vendas em quatro dos cinco ativos: WEGE3 (R\$ 47,41 contra R\$ 44,79), ITUB4 (R\$ 44,07 contra R\$ 40,61), VALE3 (R\$ 83,46 contra R\$ 78,85) e PETR4 (R\$ 44,77 contra R\$ 42,81). O sinal de COMPRA surgia com RSI médio de 63,17 e preço 7,82\% acima da média de 50 pregões; o de VENDA, com RSI médio de 35,07 e preço 8,64\% abaixo dela. Em WEGE3, único ativo em alta no período, o Buy \& Hold ganhou R\$ 603,40 e a Bollinger 20/2, R\$ 5.454,99, enquanto a v1 perdeu R\$ 5.443,38 em três ciclos de saída e reentrada.
    \item \textbf{Veto de volatilidade bloqueando saídas.} Dos 41 vetos de risco, 40 decorreram da regra dura de volatilidade. Em 25 pares ativo-pregão havia consenso técnico de VENDA com posição aberta, concentrados em BBAS3 e ITUB4 entre 27/10 e 12/11, mas a conversão de VETADO em MANTER impediu a redução, justamente quando a volatilidade aumentava.
    \item \textbf{Concentração extrema e compras residuais.} Em 101 dos 153 pregões (66,01\%), a carteira teve apenas dois ativos com peso de pelo menos 2\%, com picos de 77,49\% em VALE3, 71,40\% em PETR4 e 66,05\% em WEGE3; no extremo oposto, houve oito compras com valor inferior a R\$ 2.000.
    \item \textbf{Ausência de diversidade no comitê.} Das 765 votações, 635 (83,01\%) foram unânimes e nenhuma ficou sem consenso. Com o mesmo modelo, o mesmo \textit{prompt} e as mesmas evidências, a variação da semente, mesmo com temperatura 1,0, não produziu diversidade de critério: na prática, os cinco analistas atuaram como um único classificador de tendência.
\end{enumerate}

Os achados foram agrupados em quatro gargalos: trava de caixa (achado 1), concentração e ordens residuais (achado 4), veto simétrico (achado 3) e comitê homogêneo e pró-cíclico (achados 2 e 5). O quinto achado é coerente com a literatura discutida no Capítulo~\ref{cha:relacionados}, segundo a qual o benefício da agregação de respostas de modelos depende da diversidade efetiva entre os membros, e não apenas de seu número \citep{Wang2023, Smit2024}.

\subsection{Versão v2: correções mecânicas e arrasto de caixa}
\label{sec:res-v2}

A v2 atacou os quatro gargalos com o \textit{Trim-to-Fund}, o teto de 30\% e a ordem mínima de R\$ 2.000, o veto assimétrico (R7) e três lentes no comitê (Seção~\ref{sec:arq-evolucao}). O patrimônio final subiu R\$ 3.564,86 (+3,56 pontos percentuais), para R\$ 85.743,57, mas a carteira permaneceu em último lugar. Os efeitos sobre cada gargalo foram os seguintes:

\begin{itemize}
    \item \textbf{Trava de caixa:} duas reduções \textit{Trim-to-Fund} em PETR4 (15/09 e 15/10) e reentrada de VALE3 em 01/09;
    \item \textbf{Concentração:} nenhuma compra abaixo de R\$ 2.000, 99 pregões com três ou mais ativos relevantes e 26 ajustes do compilador (15 de teto e 11 de ordem mínima);
    \item \textbf{Veto:} 12 casos de VENDA sob veto e duas saídas executadas sob R7 (ITUB4 em 27/10 e BBAS3 em 09/11), que evitaram cerca de R\$ 827,34 de perda adicional até 12/11;
    \item \textbf{Comitê:} a unanimidade caiu de 83,01\% para 15,13\%, com 499 placares de consenso mínimo (65,66\%) e 85 pareceres sem consenso.
\end{itemize}

A melhora concentrou-se em um único ativo: o ganho de R\$ 3.596,27 em PETR4 supera o ganho total da versão, e os outros quatro ativos, somados, pioraram R\$ 31,41. Três limitações persistiram. Primeiro, a quebra da unanimidade não alterou o momento das decisões na mesma proporção: o perfil de execução continuou pró-cíclico, com RSI médio de 58,54 nas compras e 40,98 nas vendas. O diagnóstico atribui esse viés ao desempate pelo analista de regime nos placares 3--2, interpretação que os dados confirmam quanto à contagem dos placares e ao perfil resultante, mas não quanto ao voto de cada cadeira. Segundo, dez das 16 vendas zeraram a posição, e o preço médio de compra superou o de venda em quatro ativos (WEGE3, VALE3, ITUB4 e BBAS3); PETR4 foi o único ativo sem saída total e o único com resultado positivo.

Terceiro, surgiu um arrasto de caixa (\textit{cash drag}). As recomendações de COMPRA caíram de 273 na v1 para 106 na v2, contra 537 de MANTER, e o teto de 30\% por ativo limitou a reaplicação do caixa: o caixa médio foi de 24,69\% do patrimônio, superou 30\% em 65 pregões, atingiu 67,74\% e terminou em 45,86\%, com a carteira posicionada em apenas dois ativos (455 PETR4 e 490 WEGE3) e R\$ 39.321,44 em caixa não remunerado. Com dois ativos e teto de 30\% cada, ao menos 40\% do capital fica estruturalmente parado. A Figura~\ref{fig:replay-caixa} mostra a trajetória do caixa nas três versões.

% Gerado por tcc2/scripts/metricas_replay.py; requer \usepackage{pgfplots}.
\begin{figure}[htbp]
\centering
\begin{tikzpicture}
\begin{axis}[
  width=\textwidth, height=7.5cm, xmin=0, xmax=153,
  xtick={0,19,39,60,83,104,125,146},
  xticklabels={abr,mai,jun,jul,ago,set,out,nov},
  xlabel={Pregão (2026)}, ylabel={Caixa (\% do patrimônio)},
  ymajorgrids, grid style={gray!20}, tick label style={font=\footnotesize},
  label style={font=\small}, scaled y ticks=false,
  yticklabel style={/pgf/number format/use comma, /pgf/number format/1000 sep={.}},
  legend style={at={(0.5,-0.2)}, anchor=north, legend columns=3,
    font=\footnotesize, draw=none, /tikz/every even column/.append style={column sep=6pt}},
  legend cell align=left, ymin=0,
]
\addplot[black, thick, densely dotted] table[x=idx, y=cx_ma_v1, col sep=comma] {figures/replay-series.csv};
\addlegendentry{Multiagentes v1}
\addplot[black, thick, dashed] table[x=idx, y=cx_ma_v2, col sep=comma] {figures/replay-series.csv};
\addlegendentry{Multiagentes v2}
\addplot[black, very thick] table[x=idx, y=cx_ma_v3, col sep=comma] {figures/replay-series.csv};
\addlegendentry{Multiagentes v3}
\end{axis}
\end{tikzpicture}
\caption{Caixa em percentual do patrimônio das três versões da carteira Multiagentes no Replay Contínuo.}
\label{fig:replay-caixa}
\end{figure}


\subsection{Versão v3: especialização do comitê e governança do gestor}
\label{sec:res-v3}

A v3 manteve a topologia 5 + 1 + 1 e alterou a composição do comitê e as regras do gestor: cinco lentes especialistas distintas, travas sobre consensos mínimos em extremos, memória contábil de posição, piso de posição-núcleo de 8\%, venda escalonada limitada a 45\% da posição, regra contra recompra mais cara e recomposição contra o arrasto de caixa (Seção~\ref{sec:arq-evolucao}). O patrimônio final foi de R\$ 100.067,36 (+0,07\%), R\$ 14.323,79 acima da v2 e R\$ 17.888,65 acima da v1. A carteira ficou acima do Buy \& Hold (+R\$ 11.276,25), da SMA 50/200 (+R\$ 11.200,68), da Equal Weight 63d (+R\$ 10.198,07) e da Mínima Variância 252/63 (+R\$ 4.252,52), e abaixo da Bollinger 20/2 ($-$R\$ 4.819,69). Os efeitos observados foram os seguintes.

\paragraph{Diversidade do comitê.} A unanimidade caiu para 4,03\% dos pareceres, com 36,64\% de placares 4--1, 45,10\% de 3--2 e 14,23\% sem consenso.

\paragraph{Inversão do perfil de execução.} A relação entre sinal e estado técnico do ativo, pró-cíclica nas duas primeiras versões, inverteu-se. O RSI médio do sinal de COMPRA passou a 35,11 e o do sinal de VENDA, a 68,96; nas operações executadas, a 39,26 nas compras e 69,00 nas vendas. A diferença entre o RSI médio das compras e o das vendas executadas passou de +24,45 (v1) e +17,56 (v2) para $-29{,}74$. As compras ocorreram, em média, com preço 3,69\% abaixo da média de 50 pregões e a 0,93\% da banda inferior; as vendas, com preço 11,09\% acima dessa média e a 0,52\% da banda superior.

\paragraph{Saídas totais e preço médio.} A v3 não zerou nenhuma posição, contra nove saídas na v1 e dez na v2, e manteve os cinco ativos com peso de pelo menos 2\% em todos os 153 pregões. O preço médio de venda superou o de compra nos cinco ativos: PETR4 (R\$ 43,18 para R\$ 46,43), WEGE3 (R\$ 48,13 para R\$ 48,61), ITUB4 (R\$ 42,14 para R\$ 44,47), BBAS3 (R\$ 20,37 para R\$ 21,48) e VALE3 (R\$ 75,77 para R\$ 81,13).

\paragraph{Proventos e caixa.} Por nunca zerar posições antes das datas-com, a carteira recebeu os 20 eventos de proventos do período (R\$ 1.425,30 líquidos), contra nove eventos na v1 (R\$ 567,02) e nove na v2 (R\$ 509,28). O caixa médio caiu de 24,69\% para 14,94\% do patrimônio e o caixa final, de 45,86\% para 0,75\% (Figura~\ref{fig:replay-caixa}).

\paragraph{Resultado por ativo.} PETR4 (+R\$ 3.222,19), WEGE3 (+R\$ 1.991,89) e ITUB4 (+R\$ 1.435,82) passaram a resultado positivo; WEGE3 e ITUB4 respondem por R\$ 12.089,35 (84,4\%) do ganho sobre a v2. BBAS3 ($-$R\$ 1.997,86) e VALE3 ($-$R\$ 4.584,68), os dois ativos que mais caíram, permaneceram negativos. Em VALE3, a v3 teve o pior resultado das três versões e ficou abaixo do Buy \& Hold ($-$R\$ 4.255,30): o piso de posição-núcleo, que evita saídas totais, também impõe exposição a ativos em queda persistente.

\paragraph{Mecanismos efetivamente acionados.} Dos 80 ajustes registrados pelo compilador, 56 foram de piso núcleo e venda escalonada (\texttt{CORE\_FLOOR\_STEP\_SELL}), 12 de teto de peso, 9 de ordem mínima, 2 de caixa insuficiente e 1 de teto da recomposição; houve duas compras de recomposição. Não foram acionados, no período, as travas de extremo do comitê, o \textit{Trim-on-Strength}, o \textit{Trim-to-Fund}, o corte por recompra mais cara nem a saída sob R7. As 52 vendas ocorreram todas sob consenso técnico de VENDA com risco aprovado, dosadas pelo piso núcleo. O efeito observado deve-se, portanto, à combinação de lentes, piso núcleo, venda escalonada, teto, ordem mínima e recomposição, e não aos mecanismos que permaneceram inativos. Os 80 ajustes indicam ainda que a proposta bruta do gestor violava as regras com frequência, sobretudo ao pedir vendas maiores que as permitidas: a governança determinística moldou o resultado tanto quanto os sinais.

\paragraph{Efeitos colaterais.} Das 52 vendas, 27 tiveram valor inferior a R\$ 2.000 (22 abaixo de R\$ 500), somando R\$ 7.970,68: são vendas residuais produzidas pelo piso núcleo, que, com a posição próxima de 8\% do patrimônio, libera apenas as poucas ações excedentes a cada novo sinal de VENDA. Esse padrão é favorecido pelo custo puramente proporcional do replay e seria inviável com custo fixo ou lote mínimo. Os quatro pregões sem análise da v3 ocorreram após chamadas bem-sucedidas, na validação da resposta do gestor: em três deles (01/07, 01/09 e 02/10), o gestor expressou ``não vender'' como venda de quantidade zero para ativos junto ao piso, e no quarto (01/10) a resposta não cobria o universo completo.

\section{Discussão}
\label{sec:res-discussao}

\paragraph{O que os resultados permitem afirmar.} Primeiro, no replay, H1 e H2 não são sustentadas: H0 não é rejeitada para o contraste de tendência contra pesos iguais nem para a família de contrastes da carteira Multiagentes pré-especificada. Segundo, a versão pré-especificada teve, descritivamente, o menor patrimônio final, o maior \textit{drawdown} máximo e o maior giro entre as oito carteiras, e a auditoria identificou os mecanismos que produziram esse resultado; trata-se de um resultado negativo, registrado como tal. Terceiro, dentro desta trajetória e sob estas regras, a especialização das lentes e a governança determinística mudaram de forma mensurável o comportamento da carteira: menos unanimidade, inversão do momento de compra e venda, fim das saídas totais, menor concentração, captura integral de proventos e patrimônio final R\$ 17.888,65 acima da linha de base. Quarto, a comparação entre v1, v2 e v3 indica que, em um comitê de instâncias do mesmo modelo, o número de agentes importa menos que a diferença entre as funções atribuídas a eles: sementes distintas não produziram diversidade, e lentes duplicadas formaram blocos de voto que reproduziram o viés do bloco majoritário.

\paragraph{O que os resultados não permitem afirmar.} Os resultados não permitem afirmar que a v3 seja superior às referências fora deste período, que a ordenação se repita em outro regime, que algum mecanismo isolado seja responsável pelo ganho ou que o resultado represente benefício financeiro esperado. As razões são as ameaças da Seção~\ref{sec:met-ameacas}, que se manifestam concretamente nos resultados:

\begin{itemize}
    \item \textbf{Ajuste dentro da amostra.} A v2 e a v3 foram desenhadas após a observação das campanhas anteriores nos mesmos 153 pregões, com limiares escolhidos com conhecimento do período, e a v3 foi desenhada quando já se sabia que a Bollinger 20/2 era a referência de melhor resultado. Não há amostra de validação. A melhora de 17,89 pontos percentuais da v1 para a v3 mede a capacidade da arquitetura de incorporar correções dentro da amostra, e não desempenho esperado fora dela \citep{White2000, Bailey2017}.
    \item \textbf{Intervenções em bloco.} Cada versão alterou vários fatores simultaneamente, sem ablação, e parte dos mecanismos da v3 não foi acionada. A contribuição de cada pilar não é identificável.
    \item \textbf{Possível contaminação.} As campanhas foram executadas em 13, 15 e 16/11/2026, depois do período, com um modelo cujo corte de conhecimento declarado não é verificável e cujo identificador não garante versão imutável. O gestor recebe \textit{tickers}, datas e preços. A contaminação, se existir, favorece a carteira Multiagentes, não afeta as referências e, portanto, não afeta o contraste de H1 \citep{Glasserman2023, Golchin2024}.
    \item \textbf{Uma execução por versão.} Com temperatura 1,0, cada campanha é uma realização de um processo estocástico; sem repetições, a variância entre execuções não é estimada, e as diferenças entre versões incluem esse componente aleatório, além das diferenças no número de pregões sem análise (0, 1 e 4) e na data de execução sob o mesmo \textit{alias} do modelo.
    \item \textbf{A melhor referência é uma regra simples.} A Bollinger 20/2, determinística e sem custo de inferência, terminou R\$ 4.819,69 acima da v3, com Sharpe maior (0,73 contra 0,09) e \textit{drawdown} máximo menor ($-8{,}52\%$ contra $-12{,}80\%$). As lentes da v3 codificam, em boa medida, a lógica de reversão à média que caracteriza essa referência. A leitura cautelosa é que a arquitetura multiagente se mostrou capaz de expressar uma política de reversão à média com governança de carteira, e não que tenha descoberto uma política superior.
    \item \textbf{Dependência de regime.} Quatro dos cinco ativos caíram e a referência de tendência ficou próxima do Buy \& Hold. Em um período de tendência persistente, vender força e comprar fraqueza poderia inverter a ordenação das versões e das referências; o replay não permite avaliar essa possibilidade.
    \item \textbf{Custo de inferência.} Cada campanha realizou cerca de 4.703 chamadas lógicas ao modelo, com custo conhecido estimado de US\$ 7,44 (v1), US\$ 8,92 (v2) e US\$ 8,63 (v3), não deduzido do patrimônio. As referências não têm esse custo. A v3 consumiu, assim, cerca de US\$ 8,63 em inferência para terminar abaixo de uma regra determinística.
    \item \textbf{Propriedades das métricas.} Com retornos médios negativos, o Sharpe penaliza as carteiras menos voláteis, o que afeta os contrastes que envolvem a SMA 50/200; com taxa livre de risco nula, todos os índices ficam acima do que seriam sob uma taxa de juros positiva. A execução simplificada, sem custo fixo, \textit{spread} ou impacto, favorece estratégias com muitas operações pequenas, como a v3.
\end{itemize}

Em síntese, o replay oferece evidência descritiva sobre como mudanças de arquitetura alteram o comportamento de uma carteira multiagente em uma trajetória conhecida, mas não oferece evidência de superioridade dos agentes sobre as referências. O passo que poderia converter os achados exploratórios em evidência fora da amostra é a avaliação prospectiva de uma versão congelada da carteira, sem novos ajustes, no regime descrito na Seção~\ref{sec:met-protocolo}.

\section{Acompanhamento Prospectivo Preliminar}
\label{sec:obs-operacional}

Independentemente do replay, o acompanhamento diário prospectivo da carteira multiativos (Seção~\ref{sec:met-protocolo}) teve início em 14/10/2026, com as carteiras Multiagentes e Buy \& Hold montadas com a mesma cesta inicial (DP-20) e capital de R\$ 100.000,00 cada, em regime de simulação, sem capital real nem envio de ordens. O ciclo diário autônomo (DP-34), descrito na Seção~\ref{sec:arq-operacao-diaria}, foi adotado a partir de 20/10/2026. Esse acompanhamento não integra os resultados anteriores: suas carteiras, sua cesta e sua data de início são distintas das do replay, e suas decisões são registradas antes de o desfecho existir.

A série diária das quatro carteiras do acompanhamento foi exportada do servidor de operação em 17/11/2026 (\texttt{tcc2/evidence/06-operacao-diaria/serie-diaria-20261117/}) e cobre 23 pregões fechados, de 14/10 a 16/11/2026. Além de Multiagentes e Buy \& Hold, a série inclui as carteiras SMA 50/200 e Bollinger 20/2 (DP-36), criadas sobre a mesma cesta de 14/10. Nenhuma das duas regras emitiu sinal no período, e suas curvas coincidem com a do Buy \& Hold. O Buy \& Hold não negociou após a montagem da cesta; sua trajetória, de R\$ 99.835,33 em 14/10 para R\$ 82.391,17 em 16/11 ($-17{,}61\%$), reflete apenas a variação de preços dos cinco ativos. A Multiagentes realizou 19 operações após a entrada e terminou em R\$ 76.059,11 ($-23{,}94\%$). Ela esteve à frente do Buy \& Hold até 23/10, com vantagem máxima de R\$ 2.265,83 nessa data, mas caiu mais na queda de 26/10 (R\$ 84.601,02 contra R\$ 86.869,15) e ampliou a desvantagem nas semanas seguintes, até R\$ 6.332,06 em 16/11. Nas duas carteiras, o \textit{drawdown} máximo do período coincide com a perda acumulada, pois o último fechamento é também o vale.

Essas observações cobrem poucas semanas e uma queda concentrada de mercado, e não sustentam métricas de risco nem inferência sobre as hipóteses. Seu papel no trabalho é registrar o início do regime de avaliação compatível com a DP-31, no qual a contaminação do modelo pelo resultado é excluída por construção.
